% ===========================================================================
%  sviluppatore/01-introduzione/02-mappa.tex
% ===========================================================================
\section{Le due metà}\label{sez:due-meta}

La prima cosa da sapere di \Lyra{} è che non è un'applicazione: sono due, che
girano in posti diversi e si parlano attraverso tre soli canali.

\figuraTikz{mappa-sistema}
  {Le due metà del sistema. Tutto ciò che non attraversa uno dei tre canali
   numerati resta dalla parte in cui è nato.}
  {mappa-sistema}

\subsection{La metà che gira sul computer dell'utente}

È il progetto Python: \pkgCore, \pkgStrategie, \pkgBacktest, \pkgLive,
\pkgReport, e l'interfaccia Flask in \pkgDashboard. Gira in locale,
tipicamente su \url{http://127.0.0.1:5000}, e si impacchetta come applicazione
desktop (\file{desktop\_app.py} apre la stessa dashboard in una finestra
nativa invece che in una scheda del browser).

Qui stanno i dati di mercato e i risultati dei backtest, e qui restano. Non
c'è nessun percorso per cui un prezzo scaricato o un'equity line calcolata
finisca in rete.

\subsection{La metà che gira in rete}

È il sito, un'applicazione Next.js con un database Postgres. Tiene gli
account, le password, gli abbonamenti, i pagamenti e il catalogo del
marketplace. Non ha accesso ai dati di mercato e non esegue calcoli
quantitativi.

\subsection{I tre canali, e perché sono tre}

\begin{description}[leftmargin = 2.2em, style = nextline, itemsep = 6pt]
  \item[\textbf{1} --- identità e permessi]
    La dashboard chiede al sito «chi è questo utente, e che cosa può fare».
    La risposta è un \term{token} personale con dentro il livello di
    abbonamento e il flag di amministratore. Senza di esso l'applicazione
    funziona lo stesso: non si entra nel marketplace, e basta.
  \item[\textbf{2} --- chiavi e impostazioni]
    Le variabili d'ambiente --- chiavi Polygon, FRED, Stripe, SMTP --- sono
    gestite dal sito, non da un \file{.env} locale. È il verso opposto al
    canale 1, ed è l'unico che va da destra a sinistra.
  \item[\textbf{3} --- artefatti]
    Una strategia venduta sul marketplace viaggia come pacchetto di file. Va
    su quando si pubblica, viene giù quando si installa.
\end{description}

\begin{nota}{Dove cercare un bug che «sparisce»}
Quando un comportamento si vede in un ambiente e non nell'altro, la domanda da
farsi è quale dei tre canali è in mezzo. Un marketplace vuoto in locale e
pieno sul sito è il canale 1 (token mancante o scaduto); un download dati che
fallisce solo su una macchina è il canale 2 (chiave non arrivata); un
artefatto installato che non compare è il canale 3 più una delle trappole del
sistema a plugin (\vediCap{estendere}).
\end{nota}

\section{Le due applicazioni dentro la metà locale}\label{sez:strangler-intro}

Dentro la metà locale ci sono, a loro volta, due applicazioni web: la
dashboard Flask storica e un'API FastAPI più recente, avviate separatamente e
in ascolto su porte diverse. Non è un residuo da ripulire: è una migrazione
deliberata, il cosiddetto \term{strangler fig}, e il \vediCap{architettura} la
spiega per esteso.

\figuraTikz{strangler}
  {Le due applicazioni web della metà locale. Entrambe passano dalle stesse
   porte e usano lo stesso motore: ciò che cambia è il guscio.}
  {strangler}
